\documentclass[11pt,a4paper]{article}
\usepackage[utf8]{inputenc}
\usepackage[T1]{fontenc}
\usepackage{geometry}
\geometry{margin=1in}
\usepackage{hyperref}
\usepackage{enumitem}
\usepackage{booktabs}
\usepackage{longtable}
\usepackage{xcolor}
\usepackage{titlesec}
\usepackage{amsmath,amssymb}

\titleformat{\section}{\Large\bfseries\color{blue!70!black}}{}{0em}{}
\titleformat{\subsection}{\large\bfseries\color{blue!50!black}}{}{0em}{}
\titleformat{\subsubsection}{\normalsize\bfseries\color{blue!30!black}}{}{0em}{}

\title{\textbf{Replication Plan}\\[0.5em]
\large ScaWL: Scaling $k$-WL (Weisfeiler--Lehman) Algorithms in Memory and Performance on Shared and Distributed-Memory Systems\\[0.3em]
\normalsize OSTI ID: 2587225}
\author{Auto-generated Replication Blueprint}
\date{\today}

\begin{document}
\maketitle
\tableofcontents
\newpage

%---------------------------------------------------------------------
\section{Paper Overview}
%---------------------------------------------------------------------
This paper presents ScaWL, a scalable implementation of the $k$-dimensional Weisfeiler--Lehman ($k$-WL) graph isomorphism algorithm that runs efficiently on shared-memory and distributed-memory parallel systems. The $k$-WL algorithm iteratively refines $k$-tuple colorings to test graph isomorphism, but has high memory and computational cost ($O(n^k)$ tuples). ScaWL introduces memory-efficient data structures, communication-avoiding parallelism, and demonstrates strong/weak scaling on benchmark graph datasets.

%---------------------------------------------------------------------
\section{Detailed Workflow}
%---------------------------------------------------------------------

\subsection{Phase 1: Algorithm Implementation}
\begin{enumerate}[label=\textbf{1.\arabic*}]
  \item \textbf{Implement the standard $k$-WL algorithm.}
    \begin{itemize}
      \item Enumerate all $k$-tuples of vertices.
      \item Assign initial colors based on isomorphism type.
      \item Iteratively refine colors based on multiset of neighbor colors.
      \item Terminate when coloring stabilizes.
    \end{itemize}
  \item \textbf{Implement memory-efficient tuple representation.}
    \begin{itemize}
      \item Compressed storage of color histograms.
      \item Lazy tuple enumeration to reduce peak memory.
    \end{itemize}
\end{enumerate}

\subsection{Phase 2: Parallelization}
\begin{enumerate}[label=\textbf{2.\arabic*}]
  \item \textbf{Shared-memory parallelism} (OpenMP or threading).
    \begin{itemize}
      \item Parallelize tuple enumeration and color refinement.
      \item Thread-safe hash maps for color assignment.
    \end{itemize}
  \item \textbf{Distributed-memory parallelism} (MPI).
    \begin{itemize}
      \item Partition tuples across MPI ranks.
      \item Communication of boundary colors between ranks.
      \item Load balancing strategies.
    \end{itemize}
\end{enumerate}

\subsection{Phase 3: Dataset Preparation}
\begin{enumerate}[label=\textbf{3.\arabic*}]
  \item \textbf{Obtain benchmark graphs.}
    \begin{itemize}
      \item CFI (Cai--F\"{u}rer--Immerman) graphs for hard isomorphism instances.
      \item Random regular graphs.
      \item Real-world graphs from the TU Dortmund / benchmark collections.
    \end{itemize}
  \item \textbf{Prepare graph input formats} compatible with ScaWL.
\end{enumerate}

\subsection{Phase 4: Experiments}
\begin{enumerate}[label=\textbf{4.\arabic*}]
  \item \textbf{Correctness verification}: run $k$-WL on known isomorphic/non-isomorphic graph pairs.
  \item \textbf{Strong scaling}: fix problem size, increase core/node count; plot runtime vs.\ cores.
  \item \textbf{Weak scaling}: increase problem size proportional to core count.
  \item \textbf{Memory scaling}: measure peak memory vs.\ $k$ and $n$.
  \item \textbf{Compare against serial/naive implementations.}
\end{enumerate}

\subsection{Phase 5: Validation}
\begin{enumerate}[label=\textbf{5.\arabic*}]
  \item Reproduce scaling plots from the paper.
  \item Reproduce speedup and efficiency tables.
  \item Verify that ScaWL produces identical colorings to the serial reference.
\end{enumerate}

%---------------------------------------------------------------------
\section{Datasets}
%---------------------------------------------------------------------

\begin{longtable}{p{4.5cm} p{3.5cm} p{6cm}}
\toprule
\textbf{Dataset / Source} & \textbf{Type} & \textbf{Location / Notes} \\
\midrule
\endfirsthead
\toprule
\textbf{Dataset / Source} & \textbf{Type} & \textbf{Location / Notes} \\
\midrule
\endhead
CFI graphs & Synthetic & Generated programmatically (standard construction) \\
\midrule
Random regular graphs & Synthetic & Generated with NetworkX or custom code \\
\midrule
TU Dortmund benchmark & Public & \url{https://chrsmrrs.github.io/datasets/} \\
\midrule
ScaWL repository test graphs & Bundled & From the paper's code repository \\
\bottomrule
\end{longtable}

%---------------------------------------------------------------------
\section{Models, Tools, and Software}
%---------------------------------------------------------------------

\begin{longtable}{p{4.5cm} p{2.5cm} p{6.5cm}}
\toprule
\textbf{Tool / Library} & \textbf{Version} & \textbf{Purpose / URL} \\
\midrule
\endfirsthead
\toprule
\textbf{Tool / Library} & \textbf{Version} & \textbf{Purpose / URL} \\
\midrule
\endhead
ScaWL & latest & Paper's implementation \newline GitHub (see paper for URL) \\
\midrule
C/C++ compiler & GCC $\geq$ 10 or Clang $\geq$ 12 & Compilation of ScaWL \\
\midrule
MPI (OpenMPI / MPICH) & system & Distributed-memory parallelism \\
\midrule
OpenMP & system & Shared-memory parallelism \\
\midrule
CMake & $\geq 3.18$ & Build system \\
\midrule
Python / NetworkX & $\geq 3.8$ & Graph generation and verification \\
\midrule
Matplotlib & $\geq 3.4$ & Scaling and performance plots \\
\midrule
nauty / Traces & latest & Reference graph isomorphism solver for validation \newline \url{https://pallini.di.uniroma1.it/} \\
\bottomrule
\end{longtable}

\subsection{Hardware Requirements}
\begin{itemize}
  \item HPC cluster or multi-node system for distributed-memory scaling experiments.
  \item $\geq$ 64\,GB RAM per node for large $k$-WL instances.
  \item For $k=3$ on graphs with $n > 100$: memory can exceed 100\,GB.
\end{itemize}

\end{document}
